\section{Method}
\label{sec:method}

\todo{Expand derivations; cite the paper/foundation equation numbers per CLAUDE.md.}

\subsection{Modes as native contact constraints}
\label{sec:method:constraint}
Each body carries mass-normalized linear modes $\Phimat$ with amplitudes $\amp$ (the
code's $q$). A contact with unit normal $\nvec$ acts on the \emph{deformed} contact
point of both participants, so the unilateral gap becomes
\begin{equation}
  C \;=\; C_{\mathrm{rigid}} \;+\; \big(\Ghat_A\!\cdot\amp_A - \Ghat_B\!\cdot\amp_B\big),
  \qquad \Ghat_X = \nvec^{\!\top} R_X\, \Phimat_X ,
  \label{eq:gap}
\end{equation}
where $R_X$ is body $X$'s rotation. Writing $z=(v,\omega)$ for a body's rigid
velocity DOFs, the full gradient of one contact row (lever arms $r_A,r_B$ from each
COM) is the ordinary rigid Jacobian extended by one modal column block per
participant:
\begin{equation}
  \frac{\partial C}{\partial z_A} = \begin{bmatrix}\nvec\\ r_A\!\times\!\nvec\end{bmatrix}\!,
  \;\;
  \frac{\partial C}{\partial z_B} = -\begin{bmatrix}\nvec\\ r_B\!\times\!\nvec\end{bmatrix}\!,
  \;\;
  \frac{\partial C}{\partial \amp_A} = +\Ghat_A,
  \;\;
  \frac{\partial C}{\partial \amp_B} = -\Ghat_B.
  \label{eq:grad}
\end{equation}
A single contact multiplier $f\ge 0$ then (i)
pushes the rigid bodies apart and (ii) drives each participant's modes by
$+\Ghat_A f$ on $A$ and $-\Ghat_B f$ on $B$. Newton's third law is built in, so the
coupling is two-way and momentum-consistent by construction --- not a post-solve
kick. The support contact (body$\leftrightarrow$ground) is the \emph{same} row type
with the ground as degenerate participant: the surface is
$y_{\mathrm{rest}}+U_y^{\!\top}q_s$, so $\partial C/\partial q_s=-U_y$ with $U_y$
the support mode shapes' vertical components at the contact's footprint. That is
the whole generalization --- every body with modes participates, including stacked
bodies that only ever touch another body. Mass-normalized modes give
$M_q=I$, $K_q=\operatorname{diag}(\omega_j^2)$, $D_q=\alpha_0 I+\alpha_1 K_q$
(Rayleigh; DCR Eq.~7), so no reduced inverse is ever formed.
\note{This unified $(z,Q)$ row is documented in \texttt{two\_band\_coupling.html} and
\texttt{docs/sheldon\_report/README.md}.}

\subsection{Solving the stiff modal block in a local iteration}
\label{sec:method:solve}
A global QP absorbs a stiff modal DOF ($K_q\!\sim\!10^{11}$) implicitly for free; a
fixed low-iteration \emph{local} solver makes that stiff, dense modal block diverge
under naive Gauss--Seidel. Making it stable \emph{and} two-way there is the
algorithmic content of this section, and the reason the passivity bound of
\S\ref{sec:method:passivity} is needed. We embed the row of Eq.~\ref{eq:grad} in
both solver families.

\paragraph*{AVBD (augmented Lagrangian).} Each contact carries a penalty-plus-
multiplier force $f=\max(0,\rho\,C+\lambda)$; body blocks solve their usual
$6\times6$ systems with Hessian $M/h^2+\sum\rho\,J^{\!\top}J$. The modal amplitudes
are chased in a separate \emph{engaged-gated block Gauss--Seidel} pass after each
body sweep: only rows with an active multiplier contribute, and the modal block is
solved implicitly (backward Euler) as
\begin{equation}
  \Big(\tfrac{1}{h^2}M_q+\tfrac{1}{h}D_q+K_q\Big)\,\Delta\amp
  \;=\; \textstyle\sum_{\mathrm{engaged}}\Ghat^{\!\top} f
        \;-\; K_q\,\amp \;-\; \tfrac{1}{h}D_q\,(\amp-\amp^{\mathrm{prev}}),
  \qquad
  \amp \leftarrow \amp + \omega_r\,\Delta\amp,
  \label{eq:qblock}
\end{equation}
with under-relaxation $\omega_r$ (\texttt{modal\_relax}, $0.7$ in production; the
diagonal $M_q,K_q,D_q$ make the solve $r$ independent scalar updates). The
alternative --- condensing the modal columns into the contact solve by a cross-term
Schur complement --- was measured to inject energy because the eliminated
$\Delta z$ is never back-substituted at a truncated iteration count; the gated
block-GS with explicit $\Delta\amp$ is what remains stable at $8\!-\!16$
iterations.

\paragraph*{XPBD (compliant position-level).} The same row becomes a compliant
constraint with $\tilde\alpha=\alpha/h^2$ and damping $\gamma$; the generalized
inverse mass simply gains the modal terms of Eq.~\ref{eq:grad},
\begin{equation}
  w \;=\; J M^{-1} J^{\!\top}
        \;+\; \Ghat_A M_{q,A}^{-1}\Ghat_A^{\!\top}
        \;+\; \Ghat_B M_{q,B}^{-1}\Ghat_B^{\!\top},
  \qquad
  \Delta\lambda \;=\; \frac{-C-\tilde\alpha\lambda-\gamma\,\dot C}
                            {(1+\gamma)\,w+\tilde\alpha},
  \label{eq:xpbd}
\end{equation}
after which positions and amplitudes update along the gradient,
$\Delta\amp_A=+M_{q,A}^{-1}\Ghat_A^{\!\top}\Delta\lambda$ (mass-normalized:
$M_q^{-1}=I$), scaled by the same under-relaxation. The modal internal dynamics
enter as per-mode compliant rows with $\alpha_j=1/\omega_j^2$. The formulation
difference matters: the AVBD chase tolerates $\omega_r=0.7$, while the XPBD
Gauss--Seidel has a measured ceiling of $\omega_r=0.25$ on the stiff scenes
($0.5$--$0.6$ ejects bodies, ${\ge}0.65$ NaN) --- one direct motivation for the
formulation-level passivity analysis below.

\subsection{The passivity bound}
\label{sec:method:passivity}
Because the local solve under-converges, the contact can inject energy into the
modes. We cap the per-step contact$\to$modal transfer by
\begin{equation}
  \dEmod \;\le\; \eta\,\dErig , \qquad \eta\in[0,1],
  \label{eq:passivity}
\end{equation}
enforced globally across all contacts each rigid step (foundation \S15). A passive
scaling coefficient $\alpha\in[0,1]$ solves the per-step quadratic
$\dEmod(\alpha)=\alpha b + \tfrac12\alpha^2 a \le E_{\max}$ (foundation \S6),
so the injected energy never exceeds the fraction $\eta$ of the rigid kinetic energy
lost during contact resolution.
\note{DEVIATION from DCR Eq.~10 (forced IIR): the injection is a velocity-level kick
to $\dot q$ funded by rigid kinetic-energy loss, not a forced response --- see
CLAUDE.md and the foundation \S15.}
